\clearpage
\section{Verification and Testing}

The repository includes JUnit tests under \texttt{src/test/java/pcd/poool}. They exercise the game model, the physics engines, the runtime helpers, the concurrent runners, the view layer, and the benchmark utilities.

\begin{itemize}
    \item \textbf{JUnit tests}: under \texttt{src/test/java/pcd/poool};
    \item \textbf{Coverage}: model, physics engines, runtime helpers, concurrent runners, view layer, benchmark utilities.
\end{itemize}

\subsection{JPF model checking}

The verification area contains two \textbf{reduced JPF harnesses}:
\path{pcd.poool.verification.jpf.ThreadedMiniHarness} and
\path{pcd.poool.verification.jpf.TaskBasedMiniHarness}. They do not model the
full game; instead, they reuse \path{MinimalProtocolState} to check the
\textbf{essential coordination protocol}:

\begin{itemize}
    \item the board cannot be owned twice at the same time;
    \item the worker/task cannot complete before board ownership is established;
    \item pending commands are drained before publication;
    \item the final state must release ownership, publish the snapshot, and reach the finished flag.
\end{itemize}

The two JPF configuration files are \texttt{threaded-minimal.jpf} and
\texttt{taskbased-minimal.jpf}. They use the same reduced setup, but each file
has its own target harness:

\begin{verbatim}
# threaded-minimal.jpf
target=pcd.poool.verification.jpf.ThreadedMiniHarness
classpath=target/jpf-classes
search.class=gov.nasa.jpf.search.DFSearch
vm.storage.class=gov.nasa.jpf.vm.JenkinsStateSet
vm.por=true
vm.max_heap=64m
search.depth_limit=40

# taskbased-minimal.jpf
target=pcd.poool.verification.jpf.TaskBasedMiniHarness
classpath=target/jpf-classes
search.class=gov.nasa.jpf.search.DFSearch
vm.storage.class=gov.nasa.jpf.vm.JenkinsStateSet
vm.por=true
vm.max_heap=64m
search.depth_limit=40
\end{verbatim}

The actual JPF runs were executed through the Docker helper:
\begin{verbatim}
python assignment-1/verification/jpf/run_jpf.py --docker --model threaded
python assignment-1/verification/jpf/run_jpf.py --docker --model taskbased
\end{verbatim}
Here, \texttt{Python} acts as a lightweight automation layer: the project was developed across heterogeneous environments, namely Windows and Linux, so the verification workflow needed a portable scripting entry point.
Both runs completed with \texttt{no errors detected}.
The verification artifacts used for this report are available in the
repository: \href{https://github.com/pcd-26/PCD-26}{pcd-26/PCD-26}.

\begin{itemize}
    \item \textbf{Threaded model}: \texttt{states: new=774, visited=1117, backtracked=1891, end=4}; \texttt{maxDepth=35}; \texttt{instructions=22411}.
    \item \textbf{Task-based model}: \texttt{states: new=249, visited=288, backtracked=537, end=3}; \texttt{maxDepth=34}; \texttt{instructions=16409}.
\end{itemize}

The threaded model explores a larger state space because its worker thread stays alive across the protocol steps. The task-based model is smaller because the parallel work is represented by a short-lived thread that is created only when needed and joined immediately after. In both cases, JPF confirms that the bounded synchronization protocol reaches completion without deadlock or assertion failure, which is exactly the level of verification this report claims \cite{clarke1999,visser2003}.
